\begin{abstract}
Este documento describe de manera detallada y progresiva la arquitectura de ejecución de ROS~2 (Robot Operating System 2), siguiendo el flujo completo de información desde el comando escrito en la terminal hasta la instanciación de un módulo Python dentro de un nodo. Se explican los mecanismos internos de cada capa: el entorno de ejecución del sistema operativo, el rol del ejecutable \texttt{ros2} como script Python, el sistema de launch files y sus mecanismos de sustitución diferida, el ciclo de vida de arranque de un nodo, el Parameter Server, y el loop de eventos en tiempo real basado en DDS. El objetivo es que un lector con conocimientos de Python y Linux pero sin experiencia previa en ROS~2 pueda comprender no solo \emph{qué} ocurre, sino \emph{por qué} ocurre, y sea capaz de trazar el recorrido de cualquier argumento desde la terminal hasta su uso en el código.
\end{abstract}
